\section{Nghiên cứu liên quan}
\subsection{Phát hiện ảnh từ generator chưa gặp}
Wang và cộng sự huấn luyện detector trên ProGAN rồi kiểm thử trên nhiều hệ thống sinh ảnh CNN. Kết quả cho thấy preprocessing và augmentation ảnh hưởng đến khả năng chuyển giao \cite{Wang_2020_CVPR}. Ojha và cộng sự khảo sát nearest-neighbor và linear probing trên đặc trưng thị giác và ngôn ngữ pretrained để phát hiện ảnh từ các họ generator mới \cite{Ojha_2023_CVPR}.

Các công trình này là cơ sở cho cách đánh giá theo generator. FaceTrace dùng DINOv3 đóng băng và MLP trên ảnh khuôn mặt. Điểm số không được so sánh trực tiếp với các công trình trên vì dữ liệu và cách đánh giá khác nhau.

\subsection{Thông tin trong phổ ảnh}
Frank và cộng sự dùng đặc trưng tần số để nhận diện ảnh giả từ GAN \cite{Frank_2020_ICML}. Durall và cộng sự phân tích khác biệt phân bố phổ liên quan đến upsampling trong mô hình sinh ảnh \cite{Durall_2020_CVPR}.

FaceTrace dùng log-magnitude FFT theo từng kênh RGB và lấy trung bình theo bán kính. Cách biểu diễn này gọn nhưng bỏ thông tin hướng, vị trí và phase. Nhánh Forensic và các fusion được đánh giá riêng trên ảnh sạch và ảnh sau xử lý để kiểm tra giới hạn này.

\subsection{Kết hợp các loại đặc trưng}
AIDE kết hợp đặc trưng ngữ nghĩa CLIP với đặc trưng mức thấp từ các patch được chọn theo tần số \cite{AIDE_2025_ICLR}. Công trình này gợi ý việc kiểm tra thông tin bổ sung giữa các loại đặc trưng. FaceTrace dùng trung bình toàn bộ patch DINOv3 và phổ radial FFT, sau đó fusion ở xác suất. AIDE chưa được triển khai để so sánh trực tiếp.

DINOv3 cung cấp cả đặc trưng toàn ảnh và patch \cite{DINOv3_2025}. Hai loại token được lấy từ cùng backbone để hạn chế khác biệt do kiến trúc. Tên Local chỉ nói đến việc dùng patch tokens, không có nghĩa mô hình đã định vị vùng ảnh giả.

AI-Face được xây dựng cho phát hiện khuôn mặt AI và đánh giá fairness \cite{Lin_2025_CVPR}. Báo cáo này dùng một subset để kiểm tra generator chưa gặp, chưa đánh giá fairness.

\begin{table}[H]
\centering
\caption{Các công trình tham khảo khi thiết kế thí nghiệm.}
\label{tab:related-work}
\small
\begin{tabularx}{\linewidth}{p{.22\linewidth}XX}
\toprule
Công trình & Nội dung chính & Áp dụng vào FaceTrace \\
\midrule
Wang \cite{Wang_2020_CVPR} & Chuyển giao giữa hệ thống sinh CNN & Kiểm tra generator chưa gặp và xử lý ảnh \\
Ojha \cite{Ojha_2023_CVPR} & Phân loại trên đặc trưng pretrained & Baseline với backbone đóng băng \\
Frank và Durall \cite{Frank_2020_ICML,Durall_2020_CVPR} & Đặc trưng trong miền tần số & Nhánh phổ độc lập và thử ảnh sau xử lý \\
AIDE \cite{AIDE_2025_ICLR} & Kết hợp ngữ nghĩa và mức thấp & Kiểm tra lợi ích khi fusion \\
DINOv3 \cite{DINOv3_2025} & Đặc trưng tự giám sát & So sánh CLS và patch tokens \\
AI-Face \cite{Lin_2025_CVPR} & Dữ liệu khuôn mặt nhiều nguồn & Dựng tập pilot chia theo generator \\
\bottomrule
\end{tabularx}
\end{table}
